\section{Google Cloud Platform (GCP)}

Google Cloud Platform (GCP) adalah layanan komputasi awan yang dikembangkan oleh Google untuk menyediakan berbagai layanan infrastruktur, analitik data, dan kecerdasan buatan (AI). Dengan menggunakan GCP, pengguna dapat menyimpan, mengelola, dan memproses data dalam skala besar tanpa harus mengelola server fisik secara langsung. Salah satu fitur unggulan GCP adalah kemampuannya dalam mendukung pengembangan model {\it machine learning} dengan berbagai alat, seperti Google Cloud Storage untuk penyimpanan data, BigQuery untuk analitik data, dan Vertex AI untuk {\it model deployment}. Selain itu, GCP mendukung berbagai {\it framework} populer seperti {\it TensorFlow, Scikit-learn}, dan {\it PyTorch} yang memungkinkan pengguna untuk menerapkan model {\it machine learning}.

GCP memiliki beberapa keunggulan utama yang membuatnya menjadi pilihan utama dalam penerapan {\it machine learning}. Salah satu kelebihannya adalah skalabilitas dan fleksibilitas yang memungkinkan pengguna untuk menyesuaikan kapasitas penyimpanan dan daya komputasi sesuai kebutuhan. Hal ini sangat berguna ketika menangani data dalam jumlah besar atau melakukan pelatihan model dengan sumber daya komputasi yang tinggi. Selain itu, GCP menawarkan layanan {\it machine learning} terintegrasi yang mana mempermudah pengembangan, pelatihan, dan penyebaran model tanpa harus mengatur infrastruktur secara manual. Dari segi keamanan, GCP memiliki sistem keamanan tingkat tinggi, termasuk enkripsi data saat transit dan saat diam, serta kontrol akses berbasis peran untuk melindungi data pengguna. Keunggulan lainnya adalah kemampuan analitik data yang kuat di mana layanan seperti BigQuery memungkinkan analisis data dalam skala besar dengan performa tinggi. Namun, GCP juga memiliki beberapa kekurangan yang perlu diperhatikan. Biaya penggunaan GCP bisa menjadi mahal jika tidak dikonfigurasi dengan baik, terutama untuk proyek {\it machine learning} skala besar yang membutuhkan banyak sumber daya komputasi. 

Untuk mengimplementasikan model deteksi anomali seperti {\it Isolation Forest} dan {\it One-Class SVM} di GCP, dibutuhkan serangkaian tahapan mulai dari persiapan data hingga penyebaran model. Tahapan ini bersifat umum dan fleksibel untuk model-model berbasis Python dengan langkah-langkah sebagai berikut\footnote{\url{https://medium.com/publicis-sapient-france/how-to-deploy-your-own-ml-model-to-gcp-in-5-simple-steps-bf2b5898c1ab}, diakses pada 17 Maret 2025}.
    
    \begin{enumerate}
      \item {Persiapan data}
      \begin{enumerate}
        \item Mengumpulkan data menggunakan {\it data capture} yang terhubung dengan GCP. Data ini mencakup identitas penyewa, metode pembayaran, waktu {\it check-in/out}, dan lama menginap. Data disimpan di Google Cloud Storage (GCS) untuk diproses lebih lanjut.
        \item Membersihkan data menggunakan \texttt{pandas}:
        \begin{lstlisting}[language=Python]
            import pandas as pd
            # Membaca dataset dari GCS
            data = pd.read_csv("data_capture.csv")
            # Menghapus nilai kosong (NaN)
            data.dropna(inplace=True)
        \end{lstlisting}
        
        \item Melakukan normalisasi data menggunakan \texttt{StandardScaler}:
        \begin{lstlisting}[language=Python]
            from sklearn.preprocessing import StandardScaler
            scaler = StandardScaler()
            data_scaled = scaler.fit_transform(data)
        \end{lstlisting}
        
        \item Menyimpan kembali data bersih ke GCS:
        \begin{lstlisting}[language=Python]
            from google.cloud import storage
            client = storage.Client()
            bucket = client.bucket("nama-bucket")
            blob = bucket.blob("cleaned_dataset.csv")
            blob.upload_from_filename("cleaned_data.csv")
        \end{lstlisting}
      \end{enumerate}
    
      \item {Pengembangan model {\it machine learning}}
      \begin{enumerate}
        \item Mengimpor pustaka yang dibutuhkan:
        \begin{lstlisting}[language=Python]
            import numpy as np
            import pandas as pd
            from sklearn.ensemble import IsolationForest
            from sklearn.svm import OneClassSVM
            from google.cloud import storage
        \end{lstlisting}
        
        \item Membaca data dari GCS:
        \begin{lstlisting}[language=Python]
            client = storage.Client()
            bucket = client.get_bucket('nama-bucket')
            blob = bucket.blob('cleaned_dataset.csv')
            data = pd.read_csv(blob.open('r'))
        \end{lstlisting}
        
        \item Membagi data menjadi {\it training set} dan {\it testing set}:
        \begin{lstlisting}[language=Python]
            from sklearn.model_selection import train_test_split
            X_train, X_test = train_test_split(data, test_size=0.2, random_state=42)
        \end{lstlisting}
        
        \item Melatih model {\it Isolation Forest} dan {\it One-Class SVM}:
        \begin{lstlisting}[language=Python]
            # Isolation Forest
            model_if = IsolationForest(n_estimators=100, contamination=0.1, random_state=42)
            model_if.fit(X_train)
            
            # One-Class SVM
            model_svm = OneClassSVM(kernel='rbf', gamma='auto')
            model_svm.fit(X_train)
        \end{lstlisting}
      \end{enumerate}
    
      \item {Evaluasi model}
      \begin{enumerate}
        \item Melakukan prediksi pada {\it testing set}:
        \begin{lstlisting}[language=Python]
            y_pred_if = model_if.predict(X_test)
            y_pred_svm = model_svm.predict(X_test)
        \end{lstlisting}
        
        \item Mengukur performa model dengan \texttt{classification\_report}:
        \begin{lstlisting}[language=Python]
            from sklearn.metrics import classification_report
            print(classification_report(y_test, y_pred_if))
            print(classification_report(y_test, y_pred_svm))
        \end{lstlisting}
      \end{enumerate}
    
      \item {Deploy model ke Google AI Platform}
      \begin{enumerate}
        \item Menyimpan model ke dalam format \texttt{.pkl}:
        \begin{lstlisting}[language=Python]
            import joblib
            joblib.dump(model_if, 'isolation_forest.pkl')
            joblib.dump(model_svm, 'one_class_svm.pkl')
        \end{lstlisting}
        
        \item Mengunggah model ke GCS:
        \begin{lstlisting}[language=Python]
            bucket.blob('models/isolation_forest.pkl').upload_from_filename('isolation_forest.pkl')
            bucket.blob('models/one_class_svm.pkl').upload_from_filename('one_class_svm.pkl')
        \end{lstlisting}
        
        \item Melakukan deploy ke Vertex AI agar dapat digunakan untuk inferensi {\it real-time}. Sistem ini akan mendukung deteksi anomali secara otomatis dan pengambilan keputusan berbasis data.
      \end{enumerate}
\end{enumerate}